\documentclass[a4paper,12pt]{article}

\usepackage[T2A]{fontenc}
\usepackage[utf8]{inputenc}
\usepackage[russian]{babel}
\usepackage{amsmath,amssymb,mathtools}
\usepackage{helvet}
\renewcommand{\familydefault}{\sfdefault}
\usepackage{icomma}
\usepackage[a4paper,left=20mm,right=20mm,top=20mm,bottom=22mm]{geometry}
\usepackage{microtype}
\usepackage{tikz}
\usetikzlibrary{calc,arrows.meta,positioning,shapes.geometric}
\usepackage{tabularx,booktabs}
\usepackage{enumitem}
\usepackage{parskip}
\usepackage{fancyhdr}
\usepackage{setspace}
\usepackage{xcolor}
\usepackage{titlesec}

\definecolor{accent}{HTML}{008080}
\definecolor{darkaccent}{HTML}{004D4D}
\definecolor{textgray}{HTML}{333333}

\color{textgray}
\onehalfspacing
\setlength{\parindent}{0pt}
\setlength{\emergencystretch}{2em}
\setlength{\headheight}{14pt}
\renewcommand{\arraystretch}{1.35}

\titleformat{\section}
  {\Large\bfseries\color{darkaccent}}
  {\thesection.}{0.6em}{}
\titleformat{\subsection}
  {\large\bfseries\color{accent}}
  {\thesubsection.}{0.5em}{}
\titlespacing*{\section}{0pt}{1.2em}{0.6em}
\titlespacing*{\subsection}{0pt}{1em}{0.4em}

\pagestyle{fancy}
\fancyhf{}
\fancyhead[L]{\small Ярослав Верещагин \textbullet\ ОГЭ}
\fancyhead[R]{\small Степени, корни, ФСУ и модуль}
\fancyfoot[C]{\thepage}
\renewcommand{\headrulewidth}{0.4pt}
\renewcommand{\footrulewidth}{0pt}

\setlist[itemize]{leftmargin=1.6em,itemsep=0.3em,topsep=0.3em}
\setlist[enumerate]{leftmargin=1.9em,itemsep=0.45em,topsep=0.3em}

\tikzset{
  flowstart/.style={
    ellipse,
    draw=darkaccent,
    line width=0.9pt,
    align=center,
    minimum width=3.8cm,
    minimum height=1.05cm,
    text width=3.4cm,
    inner sep=4pt
  },
  flowproc/.style={
    rectangle,
    rounded corners=3pt,
    draw=darkaccent,
    line width=0.9pt,
    align=center,
    minimum height=1.0cm,
    text width=7.4cm,
    inner sep=5pt
  },
  flowcond/.style={
    diamond,
    aspect=2.5,
    draw=darkaccent,
    line width=0.9pt,
    align=center,
    text width=4.0cm,
    inner sep=1.5pt
  },
  flowarrow/.style={
    -{Stealth[length=3mm,width=2mm]},
    draw=darkaccent,
    line width=0.9pt
  },
  flowlabel/.style={
    font=\small,
    fill=white,
    inner sep=1.5pt,
    text=darkaccent
  },
  flowprocsmall/.style={
    rectangle,
    rounded corners=3pt,
    draw=darkaccent,
    line width=0.9pt,
    align=center,
    minimum height=1.0cm,
    text width=3.9cm,
    inner sep=4pt
  },
  flowcondsmall/.style={
    diamond,
    aspect=1.8,
    draw=darkaccent,
    line width=0.9pt,
    align=center,
    text width=2.6cm,
    inner sep=1.2pt
  }
}

\newcommand{\ODZ}{\textbf{ОДЗ}}
\newcommand{\R}{\mathbb{R}}
\newcommand{\checkitem}{\(\square\)}

\begin{document}

% =========================================================
% ТИТУЛЬНЫЙ ЛИСТ
% =========================================================

\begin{titlepage}
\thispagestyle{empty}
\centering
\vspace*{2.2cm}

{\Large\bfseries\color{darkaccent} Чек-лист\par}
\vspace{0.7cm}

{\huge\bfseries Степени, квадратные корни,\par}
\vspace{0.15cm}
{\huge\bfseries формулы сокращённого умножения и модуль\par}

\vspace{1.1cm}

{\large Подготовка к ОГЭ по математике\par}
\vspace{0.35cm}
{\large Ярослав Верещагин\par}

\vfill

{\large \today\par}
\vspace{0.65cm}

{\normalsize
Лёвин Артём Александрович эксклюзивно для Ярослава Верещагина\par}

\vfill

\begin{tikzpicture}[x=1cm,y=1cm]
  \draw[accent,line width=1.1pt]
    (-7.2,0) .. controls (-4.6,0.7) and (-2.3,-0.55) .. (0,0)
             .. controls (2.4,0.55) and (4.7,-0.7) .. (7.2,0);
\end{tikzpicture}

\end{titlepage}

% =========================================================
% КАРТА ЗАНЯТИЯ
% =========================================================

\section{Карта занятия}

\begin{table}[ht]
\centering
\caption{Ключевые блоки повторения}
\begin{tabularx}{\linewidth}{@{}p{0.27\linewidth}X@{}}
\toprule
\textbf{Блок} & \textbf{Что требуется уметь} \\
\midrule

Степени &
Применять свойства степеней, раскрывать степень произведения
и степень степени, работать с нулевыми и отрицательными показателями. \\

Квадратные корни &
Упрощать произведения и частные под корнем при выполнении условий,
извлекать полный квадрат, сохранять модуль
в формуле $\sqrt{U^2}=|U|$. \\

Формулы сокращённого умножения &
Узнавать полный квадрат двучлена и сворачивать квадратный трёхчлен. \\

Модуль &
Понимать модуль как расстояние до нуля,
раскрывать его по знаку подмодульного выражения. \\

Комплексное преобразование &
Соединять несколько правил в одном выражении,
начинать с ОДЗ и подставлять числа после упрощения. \\

\bottomrule
\end{tabularx}
\end{table}

% =========================================================
% СТЕПЕНИ
% =========================================================

\section{Степени: базовые строительные блоки}

\subsection{Основные свойства}

Для действительных $a,b$ и целых показателей,
когда выражения определены:

\[
a^m\cdot a^n=a^{m+n},
\qquad
\frac{a^m}{a^n}=a^{m-n}
\quad (a\ne0),
\]

\[
(a^m)^n=a^{mn},
\qquad
(ab)^n=a^n b^n.
\]

Дополнительно:

\[
a^0=1
\quad (a\ne0),
\qquad
a^{-n}=\frac{1}{a^n}
\quad (a\ne0).
\]

\textbf{Важно.}
При делении показатели вычитаются только у степеней
с одинаковым основанием.

\subsection{Пример 1: сначала упростить, затем подставить}

Найдём значение

\[
\sqrt{\frac{36a^{21}}{a^{15}}}
\quad\text{при }a=2.
\]

\ODZ: $a\ne0$.

Сначала работаем со степенями:

\[
\frac{36a^{21}}{a^{15}}
=
36a^{21-15}
=
36a^6.
\]

Теперь извлекаем корень:

\[
\sqrt{36a^6}
=
6\sqrt{(a^3)^2}
=
6|a^3|.
\]

При $a=2$:

\[
6|2^3|
=
6\cdot8
=
48.
\]

\textbf{Приём.}
Числовые значения выгодно подставлять после
алгебраического упрощения.

\subsection{Пример 2: степень степени и степень произведения}

Упростим

\[
\frac{a^{23}(b^5)^4}{(ab)^{20}}.
\]

\ODZ:
\[
a\ne0,
\qquad
b\ne0.
\]

Используем свойства степеней:

\[
\frac{a^{23}(b^5)^4}{(ab)^{20}}
=
\frac{a^{23}b^{20}}{a^{20}b^{20}}
=
a^{23-20}b^{20-20}
=
a^3.
\]

Исходные ограничения $a\ne0$ и $b\ne0$
сохраняются после сокращения.

\subsection{Чётная степень и отрицательное основание}

При чётном показателе:

\[
(-a)^4=a^4.
\]

Например,

\[
\sqrt{a^8(-a)^4}
=
\sqrt{a^{12}}
=
\sqrt{(a^6)^2}
=
|a^6|.
\]

Так как

\[
a^6\ge0,
\]

получаем

\[
|a^6|=a^6.
\]

Следовательно,

\[
\sqrt{a^8(-a)^4}=a^6.
\]

При $a=2$:

\[
2^6=64.
\]

\subsection{Отрицательные показатели}

Для $a\ne0$:

\[
\frac{a^{-12}}{a^{-17}}
=
a^{-12-(-17)}
=
a^5.
\]

При $a=2$:

\[
2^5=32.
\]

\textbf{Зона особого внимания:}
двойной минус в разности показателей.

% =========================================================
% КОРНИ
% =========================================================

\section{Квадратные корни}

\subsection{Правила и условия применимости}

Для действительных выражений:

\[
\sqrt{uv}
=
\sqrt{u}\,\sqrt{v}
\quad
\text{при }u\ge0,\ v\ge0,
\]

\[
\sqrt{\frac{u}{v}}
=
\frac{\sqrt{u}}{\sqrt{v}}
\quad
\text{при }u\ge0,\ v>0.
\]

Ключевая формула:

\[
\boxed{\sqrt{U^2}=|U|}.
\]

Поэтому

\[
\sqrt{x^6}
=
\sqrt{(x^3)^2}
=
|x^3|.
\]

Если дополнительно известно $x\ge0$, то

\[
|x^3|=x^3.
\]

\subsection{Корень и произведение}

Если выполнены условия существования корней, можно записать:

\[
\sqrt{16x^6y^4}
=
\sqrt{16}\,
\sqrt{x^6}\,
\sqrt{y^4}.
\]

Получаем

\[
\sqrt{16x^6y^4}
=
4|x^3|y^2.
\]

Для суммы и разности аналогичного свойства нет:

\[
\sqrt{x+y}\ne\sqrt{x}+\sqrt{y},
\]

\[
\sqrt{x-y}\ne\sqrt{x}-\sqrt{y}
\]

в общем случае.

\subsection{Частичное извлечение и объединение корней}

Рассмотрим

\[
E=
\frac{\sqrt{25a^9}\,\sqrt{16b^8}}
{\sqrt{a^5b^8}}.
\]

Для исходного выражения требуется

\[
a>0,
\qquad
b\ne0.
\]

При этих условиях множители можно объединить под корнем:

\[
E
=
\sqrt{
\frac{
25\cdot16\,a^9b^8
}{
a^5b^8
}
}.
\]

Сокращаем степени:

\[
E
=
\sqrt{400a^4}.
\]

Следовательно,

\[
E
=
20a^2.
\]

Если $a=4$, то

\[
E=20\cdot16=320.
\]

Переменная $b$ исчезла после сокращения,
однако условие

\[
b\ne0
\]

остаётся частью ОДЗ исходного выражения.

% =========================================================
% БЛОК-СХЕМА 1
% =========================================================

\clearpage
\thispagestyle{empty}

\begin{center}
{\Large\bfseries\color{darkaccent}
Алгоритм: упрощение выражения со степенями и корнями}
\end{center}

\vspace{0.8cm}

\begin{center}
\begin{tikzpicture}[node distance=8mm and 18mm]

\node[flowstart] (start) {Начало};

\node[flowproc, below=of start] (odz)
{Запишите ОДЗ:
знаменатели $\ne0$,
подкоренные выражения $\ge0$,
знаменатели под корнем $>0$};

\node[flowproc, below=of odz] (powers)
{Уберите внешние степени:
примените
$(U^m)^n=U^{mn}$
и
$(UV)^n=U^nV^n$};

\node[flowproc, below=of powers] (samebase)
{Соберите одинаковые основания:
при умножении сложите показатели,
при делении вычтите};

\node[flowproc, below=of samebase] (roots)
{Упростите корни:
выделите полные квадраты;
объединяйте или разделяйте корни
только при выполнении условий};

\node[flowcond, below=11mm of roots] (square)
{Появилось $\sqrt{U^2}$?};

\node[flowproc, right=25mm of square] (module)
{Замените
$\sqrt{U^2}$
на
$|U|$};

\node[flowproc, below=13mm of square] (substitute)
{Если требуется числовое значение,
подставьте данные после упрощения
и выполните арифметику};

\node[flowstart, below=of substitute] (finish)
{Ответ с учётом ОДЗ};

\draw[flowarrow] (start) -- (odz);
\draw[flowarrow] (odz) -- (powers);
\draw[flowarrow] (powers) -- (samebase);
\draw[flowarrow] (samebase) -- (roots);
\draw[flowarrow] (roots) -- (square);

\draw[flowarrow]
(square.east)
--
node[flowlabel,above]{да}
(module.west);

\draw[flowarrow]
(module.south)
|-
(substitute.east);

\draw[flowarrow]
(square.south)
--
node[flowlabel,right]{нет}
(substitute.north);

\draw[flowarrow]
(substitute)
--
(finish);

\end{tikzpicture}
\end{center}

\clearpage

% =========================================================
% ФСУ
% =========================================================

\section{Формулы сокращённого умножения: работа в обратную сторону}

\subsection{Квадрат суммы и квадрат разности}

\[
(a-b)^2
=
a^2-2ab+b^2,
\]

\[
(a+b)^2
=
a^2+2ab+b^2.
\]

Эти формулы работают и при сворачивании трёхчлена:

\[
a^2-2ab+b^2
=
(a-b)^2,
\]

\[
a^2+2ab+b^2
=
(a+b)^2.
\]

\subsection{Примеры распознавания полного квадрата}

\[
x^2-4x+4
=
x^2-2\cdot x\cdot2+2^2
=
(x-2)^2,
\]

\[
9x^2-6x+1
=
(3x)^2-2\cdot3x\cdot1+1^2
=
(3x-1)^2,
\]

\[
16k^2-8k+1
=
(4k-1)^2.
\]

\textbf{Проверка полного квадрата:}

\begin{itemize}
\item первое слагаемое является квадратом;
\item последнее слагаемое является квадратом;
\item среднее слагаемое совпадает с $\pm2ab$.
\end{itemize}

% =========================================================
% МОДУЛЬ
% =========================================================

\section{Модуль}

\subsection{Смысл и определение}

Модуль числа $x$ --- расстояние от точки $x$
до нуля на числовой прямой.

Поэтому

\[
|x|\ge0
\qquad
\text{для любого }x\in\R.
\]

Определение:

\[
|x|=
\begin{cases}
x, & x\ge0,\\
-x, & x<0.
\end{cases}
\]

Для произвольного выражения $U$:

\[
|U|=
\begin{cases}
U, & U\ge0,\\
-U, & U<0.
\end{cases}
\]

\subsection{Пример раскрытия модуля}

Рассмотрим

\[
|5-2x|.
\]

По определению:

\[
|5-2x|=
\begin{cases}
5-2x, & 5-2x\ge0,\\
-(5-2x), & 5-2x<0.
\end{cases}
\]

Решаем условия.

Первый случай:

\[
5-2x\ge0
\quad\Longrightarrow\quad
x\le\frac52.
\]

Второй случай:

\[
5-2x<0
\quad\Longrightarrow\quad
x>\frac52.
\]

Итак,

\[
|5-2x|=
\begin{cases}
5-2x, & x\le\dfrac52,\\
2x-5, & x>\dfrac52.
\end{cases}
\]

\subsection{Корень из квадрата приводит к модулю}

\[
\sqrt{(2x-1)^2}
=
|2x-1|.
\]

Особенно полезна связка с формулами сокращённого умножения:

\[
\sqrt{x^2-4x+4}
=
\sqrt{(x-2)^2}
=
|x-2|.
\]

% =========================================================
% БЛОК-СХЕМА 2
% =========================================================

\clearpage
\thispagestyle{empty}

\begin{center}
{\Large\bfseries\color{darkaccent}
Алгоритм: раскрытие модуля $|U|$}
\end{center}

\vspace{1cm}

\begin{center}
\begin{tikzpicture}[node distance=11mm and 14mm]

\node[flowstart] (start2)
{Дано $|U|$};

\node[flowcondsmall, below=of start2] (known)
{Знак $U$ известен?};

\node[flowcondsmall, below=15mm of known] (nonneg)
{$U\ge0$?};

\node[flowprocsmall, left=8mm of nonneg] (keep)
{Запишите $|U|=U$};

\node[flowprocsmall, right=8mm of nonneg] (minus)
{Запишите $|U|=-U$};

\node[flowprocsmall, left=8mm of known] (piecewise)
{Знак зависит от переменной:
запишите случаи
$U\ge0$ и $U<0$};

\node[flowstart, below=34mm of nonneg] (finish2)
{Итоговое выражение};

\draw[flowarrow]
(start2)
--
(known);

\draw[flowarrow]
(known.south)
--
node[flowlabel,right]{да}
(nonneg.north);

\draw[flowarrow]
(known.west)
--
node[flowlabel,above]{нет}
(piecewise.east);

\draw[flowarrow]
(nonneg.west)
--
node[flowlabel,above]{да}
(keep.east);

\draw[flowarrow]
(nonneg.east)
--
node[flowlabel,above]{нет}
(minus.west);

\draw[flowarrow]
(keep.south)
|-
(finish2.west);

\draw[flowarrow]
(minus.south)
|-
(finish2.east);

\draw[flowarrow]
(piecewise.west)
--
++(-5mm,0)
|-
(finish2.west);

\end{tikzpicture}
\end{center}

\clearpage

% =========================================================
% КОМПЛЕКСНЫЕ ПРЕОБРАЗОВАНИЯ
% =========================================================

\section{Комплексные преобразования}

\subsection{Связка «ФСУ $\rightarrow$ квадрат $\rightarrow$ модуль»}

Рассмотрим

\[
\sqrt{a^2+8ab+16b^2}.
\]

Подкоренное выражение является полным квадратом:

\[
a^2+8ab+16b^2
=
(a+4b)^2.
\]

Следовательно,

\[
\sqrt{a^2+8ab+16b^2}
=
|a+4b|.
\]

При

\[
a=3{,}3,
\qquad
b=1{,}7
\]

получаем

\[
|3{,}3+4\cdot1{,}7|
=
|3{,}3+6{,}8|
=
|10{,}1|
=
10{,}1.
\]

\subsection{Связка «степени $\rightarrow$ корень»}

Если под корнем находится произведение степеней
одного основания, сначала удобно собрать показатель:

\[
\sqrt{a^8\cdot a^4}
=
\sqrt{a^{12}}.
\]

Далее:

\[
\sqrt{a^{12}}
=
\sqrt{(a^6)^2}
=
|a^6|.
\]

Поскольку

\[
a^6\ge0,
\]

получаем

\[
\sqrt{a^{12}}
=
a^6.
\]

\subsection{Схема самопроверки}

Перед записью ответа убедитесь:

\begin{itemize}

\item
\checkitem\
ОДЗ исходного выражения записана
и сохранена после сокращений;

\item
\checkitem\
у одинаковых оснований показатели
сложены или вычтены корректно;

\item
\checkitem\
отрицательный показатель обработан
как обратная степень;

\item
\checkitem\
при $\sqrt{U^2}$ записан модуль $|U|$;

\item
\checkitem\
корень не распределён по сумме или разности;

\item
\checkitem\
полный квадрат проверен по всем трём слагаемым;

\item
\checkitem\
числовая подстановка выполнена
после основного упрощения.

\end{itemize}

% =========================================================
% ТИПИЧНЫЕ ОШИБКИ
% =========================================================

\section{Типичные ошибки}

\begin{enumerate}

\item
\textbf{Потеря модуля}

\[
\sqrt{x^2}=|x|.
\]

\item
\textbf{Неверное распределение корня}

\[
\sqrt{x+y}\ne\sqrt{x}+\sqrt{y}
\]

в общем случае.

\item
\textbf{Ошибка при отрицательных показателях}

\[
a^{-12}:a^{-17}
=
a^{-12-(-17)}
=
a^5.
\]

\item
\textbf{Потеря ограничений после сокращения}

Если в исходном знаменателе присутствовала переменная,
условие её ненулевого значения остаётся действующим
после упрощения.

\item
\textbf{Слишком ранняя подстановка чисел}

Символьное упрощение обычно уменьшает объём вычислений
и число арифметических ошибок.

\end{enumerate}

% =========================================================
% ТРЕНИРОВОЧНЫЙ БЛОК
% =========================================================

\clearpage
\section*{Тренировочный блок}

\textbf{Указание.}
Решения оформляйте самостоятельно.
Для выражений со знаменателями и корнями
сначала фиксируйте ОДЗ.

\subsection*{Средний уровень}

\begin{enumerate}[label=\textbf{\arabic*.}]

\item
Найдите значение

\[
\sqrt{\frac{49a^{14}}{a^8}}
\quad
\text{при }a=3.
\]

\item
Упростите

\[
\frac{(x^4)^3x^2}{x^9}.
\]

\item
Представьте в виде квадрата двучлена:

\[
25t^2-20t+4.
\]

\end{enumerate}

\subsection*{Выше среднего}

\begin{enumerate}[label=\textbf{\arabic*.},resume]

\item
Найдите значение

\[
\sqrt{81x^{10}y^4}
\quad
\text{при }x=-2,\ y=3.
\]

\item
Упростите

\[
\frac{(a^7b^3)^2}{a^9b^6}.
\]

\item
Найдите значение

\[
\sqrt{36x^2-60x+25}
\quad
\text{при }x=\frac12.
\]

\item
Раскройте модуль по определению:

\[
|4-3x|.
\]

\item
Найдите значение

\[
\frac{a^{-8}}{a^{-13}}
\quad
\text{при }a=-2.
\]

\item
Найдите значение

\[
\frac{
\sqrt{9a^7}\,\sqrt{16b^6}
}{
\sqrt{a^3b^6}
}
\quad
\text{при }a=5,\ b=-2.
\]

\end{enumerate}

\subsection*{Сложная задача}

\begin{enumerate}[label=\textbf{\arabic*.},resume]

\item
Найдите значение выражения

\[
\frac{
\sqrt{(x^2-6x+9)a^{10}}
}{
\sqrt{a^4}
}
\cdot
\frac{a^{-7}}{a^{-10}}
\]

при

\[
x=1,
\qquad
a=-2.
\]

\end{enumerate}

% =========================================================
% ОТВЕТЫ
% =========================================================

\clearpage
\section*{Ответы к тренировочному блоку}

\begin{table}[ht]
\centering
\caption{Краткие ответы}

\begin{tabularx}{\linewidth}{@{}cX@{}}

\toprule
\textbf{№} & \textbf{Ответ} \\
\midrule

1 & $189$ \\

2 & $x^5$, при $x\ne0$ \\

3 & $(5t-2)^2$ \\

4 & $2592$ \\

5 & $a^5$, при $a\ne0$, $b\ne0$ \\

6 & $2$ \\

7 &
$\displaystyle
|4-3x|=
\begin{cases}
4-3x, & x\le\dfrac43,\\
3x-4, & x>\dfrac43.
\end{cases}
$ \\

8 & $-32$ \\

9 & $300$ \\

10 & $-128$ \\

\bottomrule

\end{tabularx}
\end{table}

\end{document}